% 付録A — 開発環境セットアップ
\chapter{開発環境セットアップ}\label{app:setup}

本書のコードは Python 3.12 以降と uv を使います。Stage 01、Stage 02、Stage 03 は同じ \texttt{ftqc} という import 名を持ちますが、依存関係と実装が異なる独立したプロジェクトです。混在を避けるため、それぞれのディレクトリで環境を同期して実行してください。

\section{Stage 01}

リポジトリのルートから次を実行します。
\begin{lstlisting}[language=bash]
uv sync --directory src/stages/01_mpp --locked
uv run --directory src/stages/01_mpp --locked pytest -q
uv run --directory src/stages/01_mpp --locked python examples/ch06_surface_patch.py
\end{lstlisting}

\section{Stage 02}

\begin{lstlisting}[language=bash]
uv sync --directory src/stages/02_circuit_level --locked
uv run --directory src/stages/02_circuit_level --locked pytest -q
uv run --directory src/stages/02_circuit_level --locked python examples/ch16_circuit_merge.py
\end{lstlisting}

\section{Stage 03}

Stage 03 はリポジトリ直下の \texttt{pyproject.toml} と \texttt{uv.lock} を使います。
\begin{lstlisting}[language=bash]
uv sync --locked
uv run --locked pytest -q tests
uv run --locked python examples/ch25_d3_benchmark.py
\end{lstlisting}

書籍 PDF を生成する場合は LuaLaTeX と latexmk が必要です。\texttt{book/} で \texttt{latexmk -lualatex main.tex} を実行します。依存パッケージの導入方法は利用する OS の公式手順に従ってください。
